% Math 51 self-study: Chapter 1 exercises (Section 1.8)
% Textbook: "Linear Algebra, Multivariable Calculus, and Modern Applications"
% (Stanford Math 51, version of March 25, 2021).
% The exercise statements below are paraphrased; check each against the book.
%
% Compile with:  latexmk -pdf chapter01.tex
%
% ---------- Example: how to write a solution ----------
% Each solution starts with the final answer inside an answer box,
% then gives the explanation as ordinary text after the box.
% (This example is a comment, so it does not appear in the PDF.)
% An answer box can hold several lines of sentences and displayed equations
% written with \[...\]. Two limits come from the varwidth package, which
% shrinks the box to fit: a blank line starts a new line but prints no empty
% line, and align* is pushed to the right instead of centered.
%
%   \begin{solution}
%   \begin{enumerate}
%     \item
%       \begin{answer}
%         $2\begin{bmatrix} 1 \\ 3 \end{bmatrix}
%         + \begin{bmatrix} 4 \\ -1 \end{bmatrix}
%         = \begin{bmatrix} 6 \\ 5 \end{bmatrix}$
%       \end{answer}
%       Scalar multiplication multiplies each entry by 2, giving $(2, 6)$
%       (Definition 1.3.2). Vector addition then adds matching entries:
%       $2 + 4 = 6$ and $6 + (-1) = 5$ (Definition 1.3.1).
%     \item
%       \begin{answer}
%         $\norm{\vect{v}} = 5$
%       \end{answer}
%       Explanation for part (b) goes here.
%   \end{enumerate}
%   \end{solution}

\documentclass[11pt]{article}

\usepackage[margin=1in]{geometry}  % page margins
\usepackage{amsmath,amssymb}       % bmatrix, align, \mathbb, etc.
\usepackage{amsthm}                % theorem-style environments
\usepackage{enumitem}              % (a), (b), (c) lists
\usepackage{varwidth}              % boxes that shrink to fit their contents
\usepackage[most]{tcolorbox}       % boxes; [most] enables page-breaking boxes

% ---------- Notation ----------
\newcommand{\R}{\mathbb{R}}                    % the real numbers
\newcommand{\vect}[1]{\mathbf{#1}}             % bold vector: \vect{v}
\newcommand{\norm}[1]{\left\lVert #1 \right\rVert} % length: \norm{\vect{v}}

% Row spacing in bmatrix (and every other array): n times the normal height
\renewcommand{\arraystretch}{1.2}

% ---------- Exercise and solution environments ----------
\theoremstyle{definition}
\newtheorem*{exercise}{Exercise}               % numbered by hand below
% \leavevmode\par puts "Solution." on a line of its own, so the first
% answer box has the full line width to itself.
\newenvironment{solution}
  {\begin{proof}[Solution]\leavevmode\par}
  {\end{proof}}

% Answer box: white background, sharp black border.
% hbox and varwidth upper make the box only as wide as its widest line.
% The box can still hold several lines and \[...\] displayed equations;
% text wraps once it reaches the maximum width, \linewidth-28pt. The 28pt
% leaves room for the border and padding, so the box stays inside the margin.
% nobeforeafter keeps the box on the same line as the (a), (b) label.
% baseline sets how far the label's line sits above the box's bottom edge:
% the box's height (\tcb@height) minus the border (0.8pt), the top padding
% (4pt), and the height of one line of text (\ht\strutbox). This lines the
% label up with the first line of text inside the box.
% \tcb@height contains @, which LaTeX allows in a name only between
% \makeatletter and \makeatother.
% parbox=false lets blank lines inside the box start new paragraphs, each on
% a new line with no space between paragraphs.
\makeatletter
\newtcolorbox{answer}{colback=white, colframe=black,
  boxrule=0.8pt, arc=0pt, left=6pt, right=6pt, top=4pt, bottom=4pt,
  hbox, varwidth upper=\linewidth-28pt, nobeforeafter, parbox=false,
  baseline=\tcb@height-4.8pt-\ht\strutbox,
  after=\par\medskip}
\makeatother

\setlist[enumerate,1]{label=(\alph*)}

\title{Math 51, Chapter 1 Exercises\\
  \large Vectors, vector addition, and scalar multiplication (Section 1.8)}
\author{Jeff Trevino}
\date{\today}

\begin{document}
\maketitle

% =====================================================================
\section*{Exercise 1.1}
\begin{exercise}
Practice using your PDF viewer's ``back'' command to return to your place after
following a cross-reference link in the book, and try it on the book's link to
a picture of an ellipse. (Nothing to submit.)

This back command works only in a PDF downloaded to your computer, not when
reading the book in a web browser. The book suggests downloading a fresh copy
every week, because small corrections (such as typo fixes) are made over time.
\end{exercise}

% =====================================================================
\section*{Exercise 1.2}
\begin{exercise}
Write each expression as a single column vector.
\begin{enumerate}
  \item $3\begin{bmatrix} 1 \\ 0 \\ 0 \end{bmatrix}
        - 2\begin{bmatrix} 0 \\ 1 \\ 0 \end{bmatrix}
        + 5\begin{bmatrix} 0 \\ 0 \\ 1 \end{bmatrix}$
  \item $\begin{bmatrix} 4 \\ -1 \\ 2 \\ 7 \end{bmatrix}
        - 3\begin{bmatrix} 2 \\ 0 \\ -1 \\ 1 \end{bmatrix}
        + 2\begin{bmatrix} 4 \\ -5 \\ 3 \\ 1 \end{bmatrix}$
\end{enumerate}
\end{exercise}

\begin{solution}
\begin{enumerate}
  \item
       \begin{answer}
         $\begin{bmatrix} 3 \\ -2 \\ 5 \end{bmatrix}$
       \end{answer}
       Scalar multiplication multiplies each entry by the scalar (Definition 1.3.2); vector addition adds matching entries (Definition 1.3.1).
     \item
       \begin{answer}
         $\begin{bmatrix} 6 \\ -11 \\ 11 \\ 6 \end{bmatrix}$
       \end{answer}
       As in (a), multiply each vector's entries by its scalar and then add matching entries.
\end{enumerate}
\end{solution}

% =====================================================================
\section*{Exercise 1.3}
\begin{exercise}
The geometric meaning of the algebra in this exercise is explained later, in
Figure~3.1.2, by drawing a grid of parallelograms on a plane.
\begin{enumerate}
  \item Find scalars $t$ and $t'$ with
        $t\begin{bmatrix} 1 \\ 2 \end{bmatrix}
        + t'\begin{bmatrix} 2 \\ 3 \end{bmatrix}
        = \begin{bmatrix} 4 \\ 5 \end{bmatrix}$.
        (Hint: compare entries on both sides to get two equations in two unknowns.)
  \item Find scalars $t$ and $t'$ with
        $t\begin{bmatrix} 1 \\ 2 \end{bmatrix}
        + t'\begin{bmatrix} 2 \\ 3 \end{bmatrix}
        = \begin{bmatrix} -2 \\ -1 \end{bmatrix}$.
  \item For arbitrary $x$ and $y$, find $t$ and $t'$ in terms of $x$ and $y$ with
        $t\begin{bmatrix} 1 \\ 2 \end{bmatrix}
        + t'\begin{bmatrix} 2 \\ 3 \end{bmatrix}
        = \begin{bmatrix} x \\ y \end{bmatrix}$.
        Check that $(x,y) = (4,5)$ gives your answer to (a) and
        $(x,y) = (-2,-1)$ gives your answer to (b).
\end{enumerate}
\end{exercise}

\begin{solution}
\begin{enumerate}
  \item
  	 \begin{answer}
         $t = -2$ and $ t' = 3$.
       \end{answer}
       Scalar multiplication multiplies each entry by the scalar (Definition 1.3.2) and vector addition adds matching entries (Definition 1.3.1), which gives
       
       $\begin{bmatrix} t \\ 2t \end{bmatrix} +
       \begin{bmatrix} 2t' \\ 3t' \end{bmatrix} =
       \begin{bmatrix} 4 \\ 5 \end{bmatrix}$
       
       Because vectors are equal when their matching entries are equal (Definition 1.1.2), we have two equations of two unknowns,
       
       $t + 2t' = 4$
       
       $2t + 3t' = 5$
       
       and we can solve by isolating the unknowns. Subtracting 2 times the first equation from the second gives $t' = 3$; substituting into the first equation gives $t = -2$. We confirm that  $-2\begin{bmatrix} 1 \\ 2 \end{bmatrix} + 3\begin{bmatrix} 2 \\ 3 \end{bmatrix}  = \begin{bmatrix} 4 \\ 5 \end{bmatrix}$.

       
       
  \item \begin{answer}
         $t = 4$ and $ t' = -3$.
       \end{answer}
      Scalar multiplication multiplies each entry by the scalar (Definition 1.3.2) and vector addition adds matching entries (Definition 1.3.1), which gives
       
       $\begin{bmatrix} t \\ 2t \end{bmatrix} +
       \begin{bmatrix} 2t' \\ 3t' \end{bmatrix} =
       \begin{bmatrix} -2 \\ -1 \end{bmatrix}$
       
       We have two equations of two unknowns,
       
       $t + 2t' = -2$
       
       $2t + 3t' = -1$
       
       and we can solve by isolating the unknowns. Subtracting 2 times the first equation from the second gives $t' = -3$; substituting into the first equation gives $t = 4$. We confirm that  $4\begin{bmatrix} 1 \\ 2 \end{bmatrix} - 3\begin{bmatrix} 2 \\ 3 \end{bmatrix}  = \begin{bmatrix} -2 \\ -1 \end{bmatrix}$.
  \item \begin{answer}
         $t = - 3x + 2y$ and $t' = 2x - y$.
       \end{answer}
       In terms of an arbitrary $x$ and $y$, we have
       
       $t + 2t' = x$
       
       $2t + 3t' = y$
       
       and we can solve by isolating the unknowns. Subtracting twice the first equation from second gives $t' = 2x - y$. Substitution into the second equation gives $t = - 3x + 2y$.
       
       For $(x, y) = (4, 5)$, $t = -3(4) + 2(5) = -2$ and $t' = 2(4) - 5 = 3$.
       
       For $(x, y) = (-2, -1)$, $t = -3(-2) + 2(-1) = 4$ and $t' = 2(-2) - (-1) = -3$.

\end{enumerate}
\end{solution}

% =====================================================================
\section*{Exercise 1.4}
\begin{exercise}
On square grids, draw your answers to Exercise 1.3(a) and 1.3(b).
\end{exercise}

\begin{solution}
% your answer here
\end{solution}

% =====================================================================
\section*{Exercise 1.5}
\begin{exercise}
For each vector $\vect{v}$, compute and draw
$-\vect{v}$, $2\vect{v}$, $\tfrac{1}{2}\vect{v}$, and $-\tfrac{4}{3}\vect{v}$.
\begin{enumerate}
  \item $\vect{v} = \begin{bmatrix} 2 \\ 3 \end{bmatrix}$
  \item $\vect{v} = \begin{bmatrix} 1 \\ -2 \end{bmatrix}$
\end{enumerate}
\end{exercise}

\begin{solution}
\begin{enumerate}
  \item \begin{answer}
         $\mathbf{v} = \begin{bmatrix} 2 \\ 3 \end{bmatrix}$,
         $-\mathbf{v} = \begin{bmatrix} -2 \\ -3 \end{bmatrix}$,
         $2\mathbf{v} = \begin{bmatrix} 4 \\ 6 \end{bmatrix}$,
         $\frac{1}{2}\mathbf{v} = \begin{bmatrix} 1 \\ \frac{3}{2} \end{bmatrix}$, and
         $-\frac{4}{3}\mathbf{v} = \begin{bmatrix} -\frac{8}{3} \\ -4 \end{bmatrix}$.
       \end{answer}
       Scalar multiplication multiplies each entry by the scalar (Definition 1.3.2).
  \item \begin{answer}
         $\mathbf{v} = \begin{bmatrix} 1 \\ -2 \end{bmatrix}$,
         $-\mathbf{v} = \begin{bmatrix} -1 \\ 2 \end{bmatrix}$,
         $2\mathbf{v} = \begin{bmatrix} 2 \\ -4 \end{bmatrix}$,
         $\frac{1}{2}\mathbf{v} = \begin{bmatrix} \frac{1}{2} \\ -1 \end{bmatrix}$, and
         $-\frac{4}{3}\mathbf{v} = \begin{bmatrix} -\frac{4}{3} \\ \frac{8}{3} \end{bmatrix}$.
       \end{answer}
       Scalar multiplication multiplies each entry by the scalar (Definition 1.3.2).

\end{enumerate}
\end{solution}

% =====================================================================
\section*{Exercise 1.6}
\begin{exercise}
For each pair of 2-vectors $\vect{v}, \vect{w}$, compute and draw
$\vect{v}+\vect{w}$, $\vect{v}-\vect{w}$, $2\vect{w}$, $3\vect{w}$,
$\vect{v}+2\vect{w}$, and $\vect{v}+3\vect{w}$.
\begin{enumerate}
  \item $\vect{v} = \begin{bmatrix} 3 \\ 1 \end{bmatrix}$,
        $\vect{w} = \begin{bmatrix} 1 \\ 2 \end{bmatrix}$
  \item $\vect{v} = \begin{bmatrix} 3 \\ 1 \end{bmatrix}$,
        $\vect{w} = \begin{bmatrix} -1 \\ -2 \end{bmatrix}$
\end{enumerate}
\end{exercise}

\begin{solution}
\begin{enumerate}
  \item 
  	\begin{answer}
	$\mathbf{v} + \mathbf{w} = \begin{bmatrix} 4 \\ 3 \end{bmatrix}$
	
	$\mathbf{v} - \mathbf{w} = \begin{bmatrix} 2 \\ -1 \end{bmatrix}$
	
	$2\mathbf{w} = \begin{bmatrix} 2 \\ 4 \end{bmatrix}$
	
	$3\mathbf{w} = \begin{bmatrix} 3 \\ 6 \end{bmatrix}$
	
	$\mathbf{v} + 2\mathbf{w} = \begin{bmatrix} 5 \\ 5 \end{bmatrix}$
	
	$\mathbf{v} + 3\mathbf{w} = \begin{bmatrix} 6 \\ 7 \end{bmatrix}$
	\end{answer}
	
	 Scalar multiplication multiplies each entry by the scalar (Definition 1.3.2); vector addition adds matching entries (Definition 1.3.1).
	
  \item \begin{answer}
	$\mathbf{v} + \mathbf{w} = \begin{bmatrix} 2 \\ -1 \end{bmatrix}$
	
	$\mathbf{v} - \mathbf{w} = \begin{bmatrix} 4 \\ 3 \end{bmatrix}$
	
	$2\mathbf{w} = \begin{bmatrix} -2 \\ -4 \end{bmatrix}$
	
	$3\mathbf{w} = \begin{bmatrix} -3 \\ -6 \end{bmatrix}$
	
	$\mathbf{v} + 2\mathbf{w} = \begin{bmatrix} 1 \\ -3 \end{bmatrix}$
	
	$\mathbf{v} + 3\mathbf{w} = \begin{bmatrix} 0 \\ -5 \end{bmatrix}$
	\end{answer}
	
	 Scalar multiplication multiplies each entry by the scalar (Definition 1.3.2); vector addition adds matching entries (Definition 1.3.1).
\end{enumerate}
\end{solution}

% =====================================================================
\section*{Exercise 1.7}
\begin{exercise}
Example~1.3.8 claims that for 2-vectors $\vect{v}, \vect{w}$ and
$0 \le t \le 1$, the convex combination $(1-t)\vect{v} + t\vect{w}$ lies on the
segment joining the tips of $\vect{v}$ and $\vect{w}$. For each case below,
verify this directly:
\begin{itemize}
  \item find the equation of the line $L$ through the tips of $\vect{v}$ and $\vect{w}$;
  \item check that $(1-t)\vect{v} + t\vect{w}$ satisfies that equation;
  \item use $x$-coordinates to confirm the point lies between $\vect{v}$ and $\vect{w}$.
\end{itemize}
Draw the case $t = 1/4$, showing $L$, $\vect{v}$, and $\vect{w}$. Also sketch the
region of all vectors $r\vect{v} + s\vect{w}$ with $r, s \ge 0$ and $r + s \le 1$.
(Hint: $r\vect{v} + s\vect{w} = r\vect{v} + s\vect{w} + (1-r-s)\vect{0}$ is a
convex combination.)
\begin{enumerate}
  \item $\vect{v} = \begin{bmatrix} 1 \\ 0 \end{bmatrix}$,
        $\vect{w} = \begin{bmatrix} 0 \\ 2 \end{bmatrix}$
  \item $\vect{v} = \begin{bmatrix} 1 \\ 2 \end{bmatrix}$,
        $\vect{w} = \begin{bmatrix} 2 \\ -1 \end{bmatrix}$
\end{enumerate}
\end{exercise}

\begin{solution}
\begin{enumerate}
  \item 
  	\begin{answer}
	$y = -2(x - 1)$;
	
	$2t = -2(1 - t) + 2$;
	
	$0 \le t \le 1$ gives $0 \le 1 - t \le 1$, so the point's x-coordinate lies somewhere between $\vect{v}$'s x-coordinate (1) and $\vect{w}$'s x-coordinate (0).
	\end{answer}
	The slope of a line through $\begin{bmatrix} 1 \\ 0 \end{bmatrix}$ and $\begin{bmatrix} 0 \\ 2 \end{bmatrix}$ is $-2$. The point-slope form of the line is $y = -2(x - 1)$ or $y = -2x + 2$. 
	
	$(1 - t)\vect{v} + t\vect{w} = (1-t)\begin{bmatrix} 1 \\ 0 \end{bmatrix} + t\begin{bmatrix} 0 \\ 2 \end{bmatrix} = \begin{bmatrix} 1 - t \\ 0 \end{bmatrix} + \begin{bmatrix} 0 \\ 2t \end{bmatrix} = \begin{bmatrix} 1 - t \\ 2t \end{bmatrix}$.
	
	Substituting into the line's equation, $2t = -2(1 - t) + 2$. The fact that $-2(1 - t) + 2$ simplifies to $2t$ demonstrates that the line's equation holds for every value of $t$, and every point $(1 - t)\vect{v} + t\vect{w}$ is a point on line $L$. $0 \le t \le 1$ gives $0 \le 1 - t \le 1$, so the point's x-coordinate lies somewhere between $\vect{v}$'s x-coordinate (1) and $\vect{w}$'s x-coordinate (0).
	
  \item \begin{answer}
	$y + 1 = -3(x - 2)$;
	
	$-3t + 2 = -3(t + 1) + 5$;
	
	$0 \le t \le 1$ gives $1 \le t + 1 \le 2$, so the point's x-coordinate lies somewhere between $\vect{v}$'s x-coordinate (1) and $\vect{w}$'s x-coordinate (2).
	\end{answer}
	The slope of a line through $\begin{bmatrix} 1 \\ 2 \end{bmatrix}$ and $\begin{bmatrix} 2 \\ -1 \end{bmatrix}$ is $-3$. The point-slope form of the line is $y + 1 = -3(x - 2)$ or $y = -3x + 5$. 
	
	$(1 - t)\vect{v} + t\vect{w} = (1-t)\begin{bmatrix} 1 \\ 2 \end{bmatrix} + t\begin{bmatrix} 2 \\ -1 \end{bmatrix} = \begin{bmatrix} 1 - t \\ 2 - 2t \end{bmatrix} + \begin{bmatrix} 2t \\ -t \end{bmatrix} = \begin{bmatrix}t + 1\\-3t + 2 \end{bmatrix}$.
	
	Substituting into the line's equation, $-3t + 2 = -3(t + 1) + 5$. The fact that $-3(t + 1) + 5$ simplifies to $-3t + 2$ demonstrates that the line's equation holds for every value of $t$, and every point $(1 - t)\vect{v} + t\vect{w}$ is a point on line $L$. $0 \le t \le 1$ gives $1 \le t + 1 \le 2$, so the point's x-coordinate lies somewhere between $\vect{v}$'s x-coordinate (1) and $\vect{w}$'s x-coordinate (2).

\end{enumerate}
\end{solution}

% =====================================================================
\section*{Exercise 1.8}
\begin{exercise}
Eight students (student $1, \dots, 8$) take the same five courses
$C_1, \dots, C_5$ in their first quarter at Stanford. Courses $C_1$ and $C_2$ are worth
5 units, $C_3$ is worth 4 units, and $C_4$ and $C_5$ are worth 3 units
(20 units total). Let $\vect{v}_i$ be the 8-vector whose $j$th entry is
student $j$'s grade in course $C_i$, using
$\text{A+} = 4.3$, $\text{A} = 4.0$, $\text{A}- = 3.7$, and so on.
\begin{enumerate}
  \item Each course grade is weighted in the grade point average (GPA) by the
        course's units. Write a linear combination of the $\vect{v}_i$ whose
        $j$th entry is student $j$'s GPA. Check that it is a convex combination,
        and explain why it must be convex no matter what the unit values are.
  \item Student 1 earned three A+'s and two A's; student 2 earned two A+'s and
        three A's. Assign these grades to courses so that the two students
        have the same GPA.
\end{enumerate}
\end{exercise}

\begin{solution}
\begin{enumerate}
  \item \begin{answer}
  $0.25\mathbf{v}_1 + 0.25\mathbf{v}_2 + 0.2\mathbf{v}_3 + 0.15\mathbf{v}_4 + 0.15\mathbf{v}_5$  
		\end{answer}
	For $n$-vectors $\mathbf{v}_1$, $\hdots$, $\mathbf{v}_k$, a \emph{convex combination} means a linear combination $t_1\mathbf{v}_1 + \hdots + t_k\mathbf{v}_k$ for which all $t_j \ge 0$ and $t_1 + \hdots + t_k = 1$.
	
	Weighting courses by unit proportion of twenty units gives us proportions $\frac{5}{20} = \frac{1}{4}$ for courses $C_1$ and $C_2$ (5 units of 20 units), $\frac{1}{5}$ for course $C_3$ (4 units of 20), and $\frac{3}{20}$ for courses $C_4$ and $C_5$ (3 units of 20).
	
	If $\mathbf{v}_i$ is the 8-vector whose $j$th entry is student $j$'s grade in course $C_i$, the linear combination of the $\mathbf{v}_i$ whose $j$th entry is student $j$'s GPA is
	
	$0.25\mathbf{v}_1 + 0.25\mathbf{v}_2 + 0.2\mathbf{v}_3 + 0.15\mathbf{v}_4 + 0.15\mathbf{v}_5$  
	
	This is a convex combination because units are never negative (all $t_j \ge 0$) and each course $C_i$'s weight is a proportion of the total units such that all weights sum to 1 ($t_1 + ... + t_k = 1$): $0.25 + 0.25 + 0.2 + 0.15 + 0.15 = 1$.
	
  \item 
  	\begin{answer}
	$\vect{v}_1 = \begin{bmatrix}4.0 \\ 4.3 \\ \vdots \\ v_{1,8}\end{bmatrix}$
	$\vect{v}_2 = \begin{bmatrix}4.0 \\ 4.3\\ \vdots \\ v_{2,8}\end{bmatrix}$
	$\vect{v}_3 = \begin{bmatrix}4.3 \\ 4.0\\ \vdots \\ v_{3,8}\end{bmatrix}$
	$\vect{v}_4 = \begin{bmatrix}4.3 \\ 4.0\\ \vdots \\ v_{4,8}\end{bmatrix}$
	$\vect{v}_5 = \begin{bmatrix}4.3 \\ 4.0\\ \vdots \\ v_{5,8}\end{bmatrix}$
	\end{answer}
	Classes $C_1$ and $C_2$ together contribute half of the GPA, while classes $C_3$, $C_4$, and $C_5$ together contribute the other half; these two groups are evenly weighted. This means that student 1 and student 2 will earn the same GPA if one of them earns an A+ in courses $C_1$ and $C_2$ and an A in courses $C_3$, $C_4$, and $C_5$, while the other earns an A in courses $C_1$ and $C_2$ and an A+ in courses $C_3$, $C_4$, and $C_5$. The GPA is 4.15.
	
	
	\end{enumerate}
\end{solution}

% =====================================================================
\section*{Exercise 1.9}
\begin{exercise}
Compute the distance between $\vect{v}$ and $\vect{w}$.
\begin{enumerate}
  \item $\vect{v} = \begin{bmatrix} -2 \\ 2 \end{bmatrix}$,
        $\vect{w} = \begin{bmatrix} 1 \\ -2 \end{bmatrix}$
  \item $\vect{v} = \begin{bmatrix} 3 \\ -1 \\ 4 \\ 1 \end{bmatrix}$,
        $\vect{w} = \begin{bmatrix} -2 \\ 3 \\ 3 \\ -2 \end{bmatrix}$
\end{enumerate}
\end{exercise}

\begin{solution}
\begin{enumerate}
  \item 
  	\begin{answer}
		The distance between vectors $\vect{v}$ and $\vect{w}$, $\norm{\vect{v} - \vect{w}}$, is $\sqrt{3^2 + 4^2} = \sqrt{25} = 5$.
	\end{answer}
	$\norm{\begin{bmatrix}-2\\2\end{bmatrix} - \begin{bmatrix}1\\-2\end{bmatrix}} = \norm{\begin{bmatrix}-3\\4\end{bmatrix}}
	= \sqrt{(-3)^2 + 4^2} = \sqrt{25} = 5$

	
  \item 
  	\begin{answer}
		The distance between vectors $\vect{v}$ and $\vect{w}$, $\norm{\vect{v} - \vect{w}}$, is $\sqrt{5^2 + 4^2 + 1^2 + 3^2} = \sqrt{51}$.
	\end{answer}
	
		$\norm{\begin{bmatrix}3\\-1\\4\\1\end{bmatrix} - \begin{bmatrix}-2\\3\\3\\-2\end{bmatrix}} = 
		\norm{\begin{bmatrix}5\\-4\\1\\3\end{bmatrix}}
		=\sqrt{5^2 + (-4)^2 + 1^2 + 3^2} = \sqrt{51}$
		
		The distance between two $n$-vectors $\vect{x}$ and $\vect{y}$ is $\norm{\vect{x} - \vect{y}}$ (Definition 1.6.4) and can be found by applying the $n$-dimensional version of the Pythagorean theorem (Definition 1.6.1).
\end{enumerate}
\end{solution}

% =====================================================================
\section*{Exercise 1.10}
\begin{exercise}
Compute $\norm{\vect{v}}$ and the unit vector pointing in the same direction
as $\vect{v}$.
\begin{enumerate}
  \item $\vect{v} = \begin{bmatrix} 1 \\ 3 \\ -1 \end{bmatrix}$
  \item $\vect{v} = \begin{bmatrix} -8 \\ 15 \end{bmatrix}$
  \item $\vect{v} = \begin{bmatrix} 1 \\ 1 \\ \vdots \\ 1 \end{bmatrix}$,
        an $n$-vector (answer in terms of $n$)
\end{enumerate}
\end{exercise}

\begin{solution}
\begin{enumerate}
  \item 
  	\begin{answer}
	  	$\norm{\vect{v}} = \sqrt{11}$;
	
		$\frac{\vect{v}}{\norm{\vect{v}}} = \begin{bmatrix}
			\frac{1}{\sqrt{11}}\\
			\frac{3}{\sqrt{11}}\\
			\frac{-1}{\sqrt{11}}
		\end{bmatrix}$
	\end{answer}
	
	$\norm{\vect{v}} = \sqrt{1^2 + 3^2 + (-1)^2} = \sqrt{11}$ and $\frac{\vect{v}}{\norm{\vect{v}}} =
	\begin{bmatrix}
		\frac{1}{\sqrt{11}}\\
		\frac{3}{\sqrt{11}}\\
		\frac{-1}{\sqrt{11}}
	\end{bmatrix}$
	
	$\sqrt{(\frac{1}{\sqrt{11}})^2 + (\frac{3}{\sqrt{11}})^2 + (\frac{-1}{\sqrt{11}})^2} = \sqrt{\frac{1}{11} + \frac{9}{11} + \frac{1}{11}} = 1$
	
  \item
  \begin{answer}
  $\norm{\vect{v}} = 17$;
  
      $\frac{\vect{v}}{\norm{\vect{v}}} = \begin{bmatrix}
			\frac{-8}{17} \\ \frac{15}{17}
		\end{bmatrix}$
\end{answer}
		$\norm{\vect{v}} = \sqrt{(-8)^2 + 15^2} = 17$ and $\frac{\vect{v}}{\norm{\vect{v}}} = 
		\begin{bmatrix}
			\frac{-8}{17} \\ \frac{15}{17}
		\end{bmatrix}$
  \item
  	\begin{answer}
	$\norm{\vect{v}} = \sqrt{n}$;
	
	
		$\frac{\vect{v}}{\norm{\vect{v}}} = \begin{bmatrix}
			\frac{1}{\sqrt{n}}\\
			\vdots\\
			\frac{1}{\sqrt{n}}
		\end{bmatrix}$
	\end{answer}
  
  	$\norm{\vect{v}} = \sqrt{n}$ and $\frac{\vect{v}}{\norm{\vect{v}}} = 
	\begin{bmatrix}
		\frac{1}{\sqrt{n}}\\
		\vdots\\
		\frac{1}{\sqrt{n}}
	\end{bmatrix}
	$
  
  The unit vector pointing in the same direction as $\vect{v}$ can be obtained by scaling $\vect{v}$ by $\frac{1}{\norm{\vect{v}}}$ (Section 1.7).
\end{enumerate}
\end{solution}

% =====================================================================
\section*{Exercise 1.11}
\begin{exercise}
The \emph{polarization identity} states that for all $n$-vectors $\vect{v}, \vect{w}$,
\[
  \norm{\vect{v}+\vect{w}}^2 + \norm{\vect{v}-\vect{w}}^2
  = 2\norm{\vect{v}}^2 + 2\norm{\vect{w}}^2 .
\]
(Book Figure~1.8.1 shows the parallelogram in $\R^3$ formed by two 3-vectors
$\vect{v}$ and $\vect{w}$; its diagonals are the vectors $\vect{v}+\vect{w}$ and
$\vect{v}-\vect{w}$.)
\begin{enumerate}
  \item For $n = 2$, cut the parallelogram in half along a diagonal and apply the
        Law of Cosines to each resulting triangle, taking the diagonal as the side
        opposite the angle used. Do this for both diagonals to prove the identity.
  \item For general $n$, establish the polarization identity by algebraic calculation with the definition of
“length” for $n$-vectors.\end{enumerate}
\end{exercise}

\begin{solution}
\begin{enumerate}
  \item
  	\begin{answer}
	
	The law of cosines, $c^2 = a^2 + b^2 - 2ab\cos (C)$, where $a$, $b$, and $c$ are side lengths and $C$ is the angle opposite side $c$, generalizes the Pythagorean Theorem to apply to all triangles. Assuming $n = 2$ and applying this law to the two halves of the parallelogram with $\theta$ as the angle at the origin between $\vect{v}$ and $\vect{w}$ and opposite the diagonal $\vect{v} - \vect{w}$,
	
	\vspace{\baselineskip}
	$\norm{\vect{v} - \vect{w}}^2 = 
	\norm{\vect{v}}^2 + \norm{\vect{w}}^2 - 2\norm{\vect{v}}\norm{\vect{w}}\cos (\theta)$,
	
	\vspace{\baselineskip}
	and
	
	\vspace{\baselineskip}
	$\norm{\vect{v} + \vect{w}}^2 = 
	\norm{\vect{v}}^2 + \norm{\vect{w}}^2 - 2\norm{\vect{v}}\norm{\vect{w}}\cos (180^\circ - \theta)$
	because consecutive angles of a parallelogram are supplementary. 

	\vspace{\baselineskip}	
	Because $\cos (180^\circ - \theta) = -\cos \theta$, we can rewrite the second equation as
	$\norm{\vect{v} + \vect{w}}^2 = 
	\norm{\vect{v}}^2 + \norm{\vect{w}}^2 - 2\norm{\vect{v}}\norm{\vect{w}}(-\cos\theta) = \norm{\vect{v}}^2 + \norm{\vect{w}}^2 + 2\norm{\vect{v}}\norm{\vect{w}}\cos (\theta)$.
	
	\vspace{\baselineskip}
	The final terms cancel when adding the two equations, and we get the polarization identity
	$\norm{\vect{v} + \vect{w}}^2 + \norm{\vect{v} - \vect{w}}^2 = 2\norm{\vect{v}}^2 + 2\norm{\vect{w}}^2$.
  
  	\end{answer}
  
  \item
    \begin{answer}
    
   For a general $n$,  the $n$-dimensional Pythagorean Theorem (definition 1.6.1) $\norm{\vect{v}} = \sqrt{v_1^2 + \hdots + v_n^2}$ gives us the length of an $n$-vector. And because squaring cancels a square root, the square of vector $\vect{v}$'s length is the sum of its terms' squares $v_1^2 + \hdots + v_n^2$. When we take the difference between $n$-vectors $\vect{v}$ and $\vect{w}$, the result is a new $n$-vector with entries $v_i - w_i$; summing two vectors gives us an $n$-vector with entries $v_i + w_i$.
   
   \vspace{\baselineskip}
   By the definitions of vector addition and scalar multiplication by $-1$ (Definitions 1.3.1 and 1.3.2), we can write the $i$th term of the left side of the identity according to the length as $(v_i + w_i)^2 + (v_i - w_i)^2$, and then expand terms as $v_i^2 +2v_iw_i +w_i^2 + v_i^2 - 2v_iw_i + w_i^2$  to cancel terms and arrive at $2v_i^2 + 2w_i^2$. Because the sum of the squares of a vector's entries equals the square of its length, we know \texttt{•}hat

$\sum_{i=1}^{n} [(v_i + w_i)^2 + (v_i - w_i)^2] = \sum_{i=1}^{n} [2(v_i)^2 + 2(w_i)^2]$

says the same thing as

$\norm{\vect{v} + \vect{w}}^2 + \norm{\vect{v} - \vect{w}}^2 = 2\norm{\vect{v}}^2 + 2\norm{\vect{w}}^2$.
   
    
  \end{answer}
  
\end{enumerate}
\end{solution}

% =====================================================================
\section*{Exercise 1.12}
\begin{exercise}
The \emph{$K$-means algorithm}, mentioned in Example~1.6.13, is one of the
simplest algorithms in ``unsupervised machine learning.'' This exercise explains
the algorithm and asks you to run it. You do not need to read Example~1.6.13 or
know anything about computers or machine learning.

We are given data points $\vect{x}_1, \dots, \vect{x}_m \in \R^n$ and a whole
number $K$, and we want to sort the data into $K$ clusters of points that are
near one another. The sorting should follow a step-by-step procedure that can be
applied to large data sets, especially when $n > 3$, where the data cannot be
pictured. The algorithm is:
\begin{enumerate}[label=(\arabic*)]
  \item Choose any $K$ distinct points from the data and call them
        $\vect{C}_1, \dots, \vect{C}_K$.
  \item For each $\vect{x}_i$, find the $\vect{C}_j$ nearest to it (assume there
        are no ties: no $\vect{x}_i$ is equally far from two different
        $\vect{C}_j$'s). Grouping each $\vect{x}_i$ with its nearest $\vect{C}_j$
        splits the data into $K$ clusters.
  \item For each $j$, replace $\vect{C}_j$ with the average of the $\vect{x}_i$'s
        in the $j$th cluster.
  \item Go back to Step (2) with the new $\vect{C}_j$'s from Step (3), and keep
        repeating until the process stops, meaning that some pass through
        Steps (2) and (3) ends with the same $K$ clusters it started with.
\end{enumerate}
Six consumers rate Coca-Cola (first entry) and Pepsi (second entry) from 1
(dislike very much) to 5 (like very much):
\[
  \vect{x}_1 = \begin{bmatrix} 1 \\ 1 \end{bmatrix},\;
  \vect{x}_2 = \begin{bmatrix} 1 \\ 5 \end{bmatrix},\;
  \vect{x}_3 = \begin{bmatrix} 2 \\ 4 \end{bmatrix},\;
  \vect{x}_4 = \begin{bmatrix} 3 \\ 1 \end{bmatrix},\;
  \vect{x}_5 = \begin{bmatrix} 3 \\ 2 \end{bmatrix},\;
  \vect{x}_6 = \begin{bmatrix} 4 \\ 1 \end{bmatrix}.
\]
Use $K = 2$.
\begin{enumerate}
  \item Start with $\vect{C}_1 = \vect{x}_4$ and $\vect{C}_2 = \vect{x}_5$ and run
        the algorithm. You may omit the distance computations, but give each
        pass's averages exactly, as vectors with fractional entries.
  \item Repeat with $\vect{C}_1 = \vect{x}_2$ and $\vect{C}_2 = \vect{x}_3$.
        (You should end with the same two clusters as in (a), with their
        labels swapped.)
\end{enumerate}
\end{exercise}

\begin{solution}
\begin{enumerate}
  \item
  
  \begin{answer}
  $\{\vect{x}_1, \vect{x}_4, \vect{x}_5, \vect{x}_6\}$ and $\{\vect{x}_2, \vect{x}_3 \}$
  \end{answer}
  
  If we begin with $\vect{C}_1 = \vect{x}_4$ and $\vect{C}_2 = \vect{x}_5$, then passes are as follows:
  
  \textbf{pass one}
  
 The two clusters are $ \{\vect{x}_1, \vect{x}_4, \vect{x}_6\}$ and  $\{\vect{x}_2, \vect{x}_3, \vect{x}_5\}$.
  
  We update the cluster centroids as $\vect{C}_1 = (\vect{x}_1 + \vect{x}_4 + \vect{x}_6)/3
  = \begin{bmatrix}(1 + 3 + 4) / 3 \\ (1 + 1 + 1)/3\end{bmatrix} = \begin{bmatrix}\frac{8}{3} \\ 1\end{bmatrix}$
  
  and $\vect{C}_2 = (\vect{x}_2 + \vect{x}_3 +  \vect{x}_5) / 3 = \begin{bmatrix}(1 + 2 + 3)/3\\(5 + 4 + 2)/3\end{bmatrix} = \begin{bmatrix}2\\\frac{11}{3}\end{bmatrix}$.
  
    \textbf{pass two}
    
    The two clusters are $\{\vect{x}_1, \vect{x}_4, \vect{x}_5, \vect{x}_6\}$ and $\{\vect{x}_2, \vect{x}_3\}$.
 
  The algorithm continues because $\vect{x}_5$ moved clusters, and we update the cluster centroids as
  
  $\vect{C}_1 = (\vect{x}_1 + \vect{x}_4 + \vect{x}_5 + \vect{x}_6)/4
  = \begin{bmatrix}(1 + 3 + 3 + 4) / 4 \\ (1 + 1 + 2 + 1)/4\end{bmatrix} = \begin{bmatrix}\frac{11}{4} \\ \frac{5}{4}\end{bmatrix}$
  
  and $\vect{C}_2 = (\vect{x}_2 + \vect{x}_3 ) / 2 = \begin{bmatrix}(1 + 2)/2\\(5 + 4)/2\end{bmatrix} = \begin{bmatrix}\frac{3}{2}\\\frac{9}{2}\end{bmatrix}$.
 
  
  \textbf{pass three}
    
   The two clusters are still $\{\vect{x}_1, \vect{x}_4, \vect{x}_5, \vect{x}_6\}$ and $\{\vect{x}_2, \vect{x}_3\}$, and the algorithm halts because no data points moved clusters.
    
  
  \item
  
    \begin{answer}
  $\{\vect{x}_2, \vect{x}_3 \}$ and $\{\vect{x}_1, \vect{x}_4, \vect{x}_5, \vect{x}_6\}$
  \end{answer}
  
  If we begin with $\vect{C}_1 = \vect{x}_2$ and $\vect{C}_2 = \vect{x}_3$, then the passes are as follows:
  
  \textbf{pass one}
  
 The two clusters are $ \{\vect{x}_2\}$ and  $\{\vect{x}_1, \vect{x}_3, \vect{x}_4, \vect{x}_5, \vect{x}_6\}$.
 
 We update the cluster centroids as $\vect{C}_1 = (\vect{x}_2 )/1
  = \begin{bmatrix}1 \\ 5\end{bmatrix}$ and $\vect{C}_2 = (\vect{x}_1 + \vect{x}_3 + \vect{x}_4 + \vect{x}_5 + \vect{x}_6)/5
  = \begin{bmatrix}(1 + 2 + 3 + 3 + 4) / 5 \\ (1 + 4 + 1 + 2+ 1)/5\end{bmatrix} = \begin{bmatrix}\frac{13}{5} \\ \frac{9}{5}\end{bmatrix}$.
  
  \textbf{pass two}
  
  The two clusters are $ \{\vect{x}_2, \vect{x}_3\}$ and  $\{\vect{x}_1, \vect{x}_4, \vect{x}_5, \vect{x}_6\}$.   The algorithm continues because $\vect{x}_3$ moved clusters, and we update the cluster centroids as
  $\vect{C}_1 = (\vect{x}_2 + \vect{x}_3)/2
  = \begin{bmatrix}(1 + 2)/2\\(5 + 4)/2\end{bmatrix} = \begin{bmatrix}\frac{3}{2} \\ \frac{9}{2}\end{bmatrix}$ and $\vect{C}_2 = (\vect{x}_1 + \vect{x}_4 + \vect{x}_5 + \vect{x}_6)/4
  = \begin{bmatrix}(1 + 3 + 3 + 4) / 4 \\ (1 + 1 + 2 + 1)/4\end{bmatrix} = \begin{bmatrix}\frac{11}{4} \\ \frac{5}{4}\end{bmatrix}$.
  
    \textbf{pass three}
    
   The two clusters are still $\{\vect{x}_2, \vect{x}_3\}$ and $\{\vect{x}_1, \vect{x}_4, \vect{x}_5, \vect{x}_6\}$, and the algorithm halts because no data points moved clusters. These are indeed the same two clusters as in (a), but with labels swapped.
  
\end{enumerate}
\end{solution}

\end{document}
